\documentclass[11pt]{article}
\usepackage[utf8]{inputenc}
\usepackage[T1]{fontenc}
\usepackage{geometry}
\geometry{a4paper, margin=1in}

\begin{document}

\begin{center}
\Large\bfseries Endorsement Request --- Liu (arXiv 2512.04521) Re-evaluation
\end{center}

\vspace{0.3cm}

\noindent\textbf{From:} 20242100136@dcc.edu.cn \\
\textbf{To:} wangjin@ynu.edu.cn (云南大学信息学院 王津老师) \\
\textbf{Subject:} Request for arXiv cs.LG endorsement --- Cross-domain WiFi gesture recognition reproducibility audit

\vspace{0.3cm}
\hrule
\vspace{0.3cm}

\noindent 尊敬的王老师：

您好！我是滇池学院（云南大学滇池学院）的访问研究人员，目前在做一个 WiFi CSI 跨域手势识别方向的研究。最近准备向 arXiv 投递一篇论文（cs.LG primary + eess.SP cross-list），但因本人是首次在该 archive 投稿，需要一位已经在 cs.LG 有 endorsement 资格的 author 帮忙 endorse。

了解到您长期深耕联邦学习（AAAI 2025 "Data-Free Black-Box FL"、EMNLP 2023 "FedID"）和跨域泛化（ACL 2024 "Zero-Shot Cross-Domain Dialogue State Tracking"），与我们的论文核心议题高度契合，故冒昧请求您的 endorsement。

\vspace{0.2cm}

\noindent\textbf{论文核心贡献（200 字摘要）：}

\begin{quote}
\small
WiFi CSI-based cross-domain gesture recognition has attracted intensive research in the past six years, yet reported improvements rarely reproduce under controlled evaluation. We present a \textbf{paired benchmark} covering ten recent algorithmic improvements---normalization, alignment, augmentation, reweighting, self-supervised pretraining, domain adversarial training, snapshot ensemble, and novel architectures---evaluated under an 18-fold × 5-seed cross-validated protocol on the Widar3.0 benchmark. \textbf{All ten methods fail to reach significance} ($|$Cohen's $d| < 0.3$, $p > 0.28$). Variance decomposition reveals the root cause: within-seed standard deviation (\textbf{4.82\,pp}) is 4.7× the largest reported effect, guaranteeing non-reproducible ``significant'' results in single-seed studies. We then introduce a minimal signal-representation modification---expanding the BVP operator to retain the multi-Rx antenna dimension---that yields \textbf{+8.14\,pp} (Cohen's $d = 1.71$, 19/1 wins). We further test this finding on a second dataset MMFi and find it \textbf{does not generalize}---the dataset-optimal representation \textbf{flips} to multi-scale on MMFi. \textbf{No single BVP representation is universally optimal; the choice must be matched to the dataset's geometric structure.}
\end{quote}

\vspace{0.2cm}

\noindent\textbf{论文关键文件：}
\begin{itemize}
\item \texttt{arxiv\_submission/paper\_arxiv.tex}（LaTeX 源文件，约 25 页）
\item \texttt{arxiv\_submission/paper\_arxiv.pdf}（编译后 PDF）
\item \texttt{arxiv\_submission/ARXIV\_METADATA.txt}（投稿元数据）
\item 投稿后会有 arXiv submission ID，可在 https://arxiv.org/auth/endorsers 查询
\end{itemize}

\vspace{0.2cm}

\noindent\textbf{与您研究的相关性：}
\begin{itemize}
\item 您在 AAAI 2024-2025 关于联邦学习的系列工作（如 FedID）与我们的跨域泛化主题直接相关
\item 您在 ACL 2024 关于 Zero-Shot Cross-Domain 的工作，与我们的"跨域性能不鲁棒"结论形成对照
\item 论文核心是 ML benchmark + 可复现性方法论，符合 cs.LG 的核心议题
\end{itemize}

\vspace{0.2cm}

\noindent\textbf{endorsement 操作：}
\begin{enumerate}
\item 登录您的 arXiv 账号（您应有 cs.LG 的 endorsement 资格）
\item 访问 https://arxiv.org/auth/endorsers
\item 在 endorse others 部分搜索我的 arXiv submission ID（投稿后我会在邮件补充）
\item 点击 endorse 即可（无需填写详细理由）
\end{enumerate}

\vspace{0.2cm}

\noindent 如蒙支持，万分感谢。如有任何问题，欢迎邮件沟通。

\vspace{0.3cm}

\noindent 此致 \\
敬礼

\vspace{0.3cm}

\noindent [你的名字] \\
20242100136@dcc.edu.cn \\
滇池学院（云南大学滇池学院）

\end{document}
